\documentclass[11pt,a4paper]{article}
\usepackage[utf8]{inputenc}
\usepackage[T1]{fontenc}
\usepackage{geometry}
\usepackage{graphicx}
\usepackage{hyperref}
\usepackage{amsmath}
\usepackage{amssymb}
\usepackage{booktabs}
\usepackage{tabularx}
\usepackage{xcolor}
\usepackage[hidelinks]{hyperref}

\geometry{a4paper,total={170mm,257mm},left=20mm,top=20mm}

\title{\textbf{ajedrezUPV --- Documentacion del proyecto}\\[0.5cm]
\large Motor de ajedrez con inteligencia artificial}
\author{Desarrollado en la Universitat Politecnica de Valencia (UPV)\\
\texttt{perezllorenc@gmail.com}}
\date{\today}

\begin{document}

\maketitle
\tableofcontents
\newpage

% ===========================================================================
\section{Visión general}
% ===========================================================================

\textbf{ajedrezUPV} es un motor de ajedrez con inteligencia artificial,
desarrollado en el marco de la Universitat Politecnica de Valencia (UPV).
Combina algoritmos de busqueda avanzados (alpha-beta con profundizacion
iterativa, tabla de transposicion, quiescence search, ordenacion de
movimientos) con redes neuronales para la evaluacion de posiciones.

\subsection*{Caracteristicas principales}
\begin{itemize}
  \item Motor de ajedrez en C++17 con la libreria
        \href{https://disservin.github.io/chess-library/}{Disservin chess-library}
  \item \textbf{Busqueda avanzada}: Alpha-Beta con profundizacion iterativa,
        aspiration windows, tabla de transposicion (4-way set-associative),
        quiescence search con stand-pat, ordenacion de movimientos (TT move,
        MVV-LVA, killers, history heuristic)
  \item \textbf{Evaluacion heuristica}: material con mapas de calor,
        estructura de peones (aislados, doblados, pasados), seguridad del rey
        (fraccion de casillas vecinas bajo ataque enemigo), control del
        tablero, deteccion correcta de jaque mate y tablas
  \item \textbf{Redes neuronales} (C++): red grande
        ($2565 \to 2048 \to 1536 \to 1024 \to 768 \to 512 \to 256 \to 1$,
        $\sim$14M parametros) y red junior ($782 \to 512 \to 512 \to 1$)
        con inicializacion He, perdida MSE, y serializacion JSON/binaria
  \item \textbf{Entrenamiento PyTorch}: script con mini-batches, perspectiva
        del lado al mueve, objetivo $\tanh(cp/400)$, exportacion a binario
        para inferencia C++
  \item Protocolo de juego interactivo con comandos: movimientos UCI,
        \texttt{undo}, \texttt{new}, \texttt{fen}, \texttt{d}, \texttt{go depth N}
  \item 52 tests automatizados (unitarios, perft, mate-in-N, end-to-end)
\end{itemize}

% ===========================================================================
\section{Estructura del repositorio}
% ===========================================================================

\begin{center}
\small
\begin{tabularx}{\textwidth}{@{}l X@{}}
\toprule
\textbf{Ruta} & \textbf{Descripcion} \\
\midrule
\texttt{CMakeLists.txt} & Build CMake (con FetchContent para chess-library) \\
\texttt{Makefile} & Build Make \\
\texttt{play.sh} & Script para jugar facilmente \\
\texttt{requirements.txt} & Dependencias Python \\
\midrule
\texttt{include/chess/} & Cabeceras publicas del motor \\
\quad \texttt{evaluator.h} & Interfaz base de evaluacion + constantes (\texttt{eval::} namespace) \\
\quad \texttt{general\_evaluator.h} & Evaluador general con metodos comunes \\
\quad \texttt{opening\_evaluator.h} & Evaluador de apertura (mapas de calor) \\
\quad \texttt{endgame\_evaluator.h} & Evaluador de finales (rey centralizado) \\
\quad \texttt{node\_move\_optimized.h} & Nodo del arbol + SearchResult + quiescence (legacy) \\
\quad \texttt{transposition\_table.h} & Tabla de transposicion (4-way set-associative) \\
\quad \texttt{search.h} & Busqueda on-demand con TT + ID + aspiration windows \\
\midrule
\texttt{src/engine/} & Motor de ajedrez \\
\quad \texttt{alphabeta\_optimized.cpp} & Motor principal (bucle de juego, comandos) \\
\quad \texttt{node\_move\_optimized.cpp} & Arbol optimizado con alpha-beta + quiescence (legacy) \\
\quad \texttt{transposition\_table.cpp} & Implementacion de la tabla de transposicion \\
\quad \texttt{search.cpp} & Busqueda on-demand (TT + ID + quiescence + killers) \\
\quad \texttt{evaluation/} & Evaluadores heuristicos (general, apertura, final) \\
\midrule
\texttt{src/neural/} & Redes neuronales C++ \\
\quad \texttt{network\_junior.hpp} & Red junior (782$\to$512$\to$512$\to$1) \\
\quad \texttt{square\_rook\_network.cpp} & Red grande ($\sim$14M parametros) \\
\quad \texttt{simd\_network\_wip.cpp} & Prototipo SIMD (AVX2, WIP) \\
\midrule
\texttt{src/tools/} & Herramientas de datos \\
\quad \texttt{json\_extractor.cpp} & Extractor de evaluaciones JSONL \\
\quad \texttt{json\_extractor\_junior.cpp} & Extractor + entrenador inline con perspectiva \\
\midrule
\texttt{python/} & Herramientas Python \\
\quad \texttt{train.py} & Entrenamiento PyTorch (mini-batches, tanh, export binario) \\
\quad \texttt{model.py} & Analisis de datasets y visualizacion \\
\quad \texttt{network\_reference.py} & Referencia de red neuronal (Michael Nielsen) \\
\quad \texttt{mnist\_loader.py} & Cargador MNIST (referencia) \\
\midrule
\texttt{test/} & 52 tests automatizados \\
\texttt{data/dataset.csv} & Dataset de analisis \\
\texttt{lichess\_db/} & Chunks JSONL de evaluaciones de Lichess \\
\texttt{third\_party/} & json.hpp, chess.hpp \\
\texttt{docs/} & Documentacion LaTeX y PDFs \\
\texttt{LICENSE} & CC BY-NC 4.0 \\
\bottomrule
\end{tabularx}
\end{center}

% ===========================================================================
\section{Motor de busqueda}
% ===========================================================================

El motor utiliza dos implementaciones de busqueda:

\subsection*{Search (recomendada)}
Busqueda on-demand con generacion lazy de movimientos. Implementa:
\begin{itemize}
  \item Profundizacion iterativa con aspiration windows (depth $\geq$ 4)
  \item Tabla de transposicion 4-way set-associative (16 MB por defecto)
  \item Quiescence search con stand-pat y MVV-LVA
  \item Ordenacion: TT move $\to$ capturas MVV-LVA $\to$ promociones $\to$ killers $\to$ history
  \item Deteccion de jaque mate, ahogado, tablas (50 movimientos, repeticion, material insuficiente)
  \item Control de tiempo para futura integracion UCI
\end{itemize}

\subsection*{NodeMoveOptimized (legacy)}
Arbol pre-construido con arena PMR (\texttt{std::pmr::unsynchronized\_pool\_resource}).
Mismas tecnicas de busqueda pero con el arbol materializado en memoria.
Util para depuracion y benchmarks.

\subsection*{Tabla de transposicion}
\begin{itemize}
  \item Entradas: hash (64 bits) + movimiento + puntuacion + profundidad + flag + edad
  \item Buckets de 4 entradas (set-associative)
  \item Flags: \texttt{EXACT}, \texttt{LOWER} (beta-corte), \texttt{UPPER} (fail-low)
  \item Puntuaciones de mate ajustadas por ply (\texttt{scoreForTT}/\texttt{scoreFromTT})
\end{itemize}

\subsection*{Constantes de busqueda}
\begin{verbatim}
MATE_SCORE  = 30000.0f   // Puntuacion de jaque mate
INFINITY    = 50000.0f   // Infinito de la ventana
MAX_DEPTH   = 3          // Profundidad maxima (arbol legacy)
MAX_BRANCH  = 100        // Limite de hijos por nodo (legacy)
\end{verbatim}

% ===========================================================================
\section{Evaluador heuristico}
% ===========================================================================

\subsection*{Escala}
Todos los valores estan en centipeones:
\begin{verbatim}
PAWN_VALUE   = 100    KNIGHT_VALUE = 300    BISHOP_VALUE = 315
ROOK_VALUE   = 500    QUEEN_VALUE  = 900
\end{verbatim}

\subsection*{Fases del juego}
La evaluacion detecta la fase por presencia de damas:
\begin{itemize}
  \item \textbf{Apertura} (\texttt{OpeningEvaluator}): mapas de calor para piezas menores, rey en esquinas
  \item \textbf{Final} (\texttt{EndgameEvaluator}): rey centralizado, peones avanzados
\end{itemize}

\subsection*{Terminos de evaluacion}
\begin{enumerate}
  \item \textbf{Material + posicion}: suma de valores base + mapas de calor por pieza
  \item \textbf{Estructura de peones}: penalti por peones aislados ($-0.5$), doblados ($-0.5$), bonus por pasados ($+1.0$). Comprueba todas las filas hacia adelante.
  \item \textbf{Seguridad del rey}: fraccion de casillas vecinas al rey bajo control enemigo (0 = seguro, 1 = inseguro)
  \item \textbf{Control del tablero}: casillas controladas por el color consultado, pesadas por importancia central
\end{enumerate}

\subsection*{Perspectiva}
Todos los terminos devuelven puntuacion en perspectiva del \emph{lado al mueve}
(convencion negamax). La busqueda niega el signo al subir por el arbol.

% ===========================================================================
\section{Redes neuronales}
% ===========================================================================

\subsection*{Red junior (\texttt{network\_junior.hpp})}
\begin{itemize}
  \item Arquitectura: $782 \to 512 \to 512 \to 1$
  \item Inicializacion He: $\text{weights} \sim \mathcal{N}(0, \sqrt{2/N_{in}})$
  \item Activacion: ReLU (capas ocultas), lineal (salida)
  \item Perdida: MSE ($\frac{1}{2}(\text{pred} - \text{target})^2$)
  \item Optimizador: Adam ($\beta_1=0.9$, $\beta_2=0.999$, $\epsilon=10^{-8}$)
\end{itemize}

\subsection*{Red grande (\texttt{square\_rook\_network.cpp})}
\begin{itemize}
  \item Arquitectura: $2565 \to 2048 \to 1536 \to 1024 \to 768 \to 512 \to 256 \to 1$
  \item $\sim$14 millones de parametros
  \item Perdida: Huber (con delta=1.0)
  \item Serializacion a JSON con \texttt{std::async} para paralelismo
\end{itemize}

\subsection*{Encoding de entrada (782 features)}
\begin{enumerate}
  \item 768 entradas: 12 tipos de pieza $\times$ 64 casillas (one-hot)
  \item 4 entradas: derechos de enroque
  \item 1 entrada: lado al mueve (siempre 1.0, codificado desde nuestra perspectiva)
  \item 8 entradas: casilla de captura al paso (en passant)
  \item 1 entrada: reloj de 50 movimientos (normalizado)
\end{enumerate}

Cuando mueve negro, el tablero se espeja verticalmente para que ``nuestras''
piezas siempre avancen ``hacia arriba''. Esto reduce a la mitad los patrones
que la red debe aprender.

\subsection*{Entrenamiento PyTorch (\texttt{python/train.py})}
\begin{itemize}
  \item Mini-batches de 256 ejemplos
  \item Funcion objetivo: $\tanh(cp/400) \in [-1, 1]$
  \item Weight decay ($10^{-4}$) en lugar de dropout
  \item Seeds fijos para reproducibilidad
  \item Exportacion a tres formatos: \texttt{.pt} (PyTorch), \texttt{.bin} (binario float32), \texttt{.json}
\end{itemize}

% ===========================================================================
\section{Sistema de tests}
% ===========================================================================

El proyecto incluye 52 tests automatizados organizados en 6 suites:

\begin{center}
\small
\begin{tabularx}{\textwidth}{@{}l c X@{}}
\toprule
\textbf{Suite} & \textbf{Tests} & \textbf{Que verifica} \\
\midrule
\texttt{test\_node\_move} & 3 & Arbol de nodos, rebuild, alpha-beta con restauracion de tablero \\
\texttt{test\_search\_benchmark} & 1 & Benchmark de busqueda depth 1-3 \\
\texttt{test\_json\_extractor} & 3 & Parser de evaluaciones JSONL de Lichess \\
\texttt{test\_json\_extractor\_junior} & 3 & Conversion FEN a tensor de entrada (782 features) \\
\texttt{test\_phase2} & 13 & TT, busqueda on-demand, perft, mate-in-1 y mate-in-2 \\
\texttt{test\_e2e} & 29 & End-to-end: arranque, movimientos, EOF, comandos, reglas de ajedrez \\
\bottomrule
\end{tabularx}
\end{center}

\subsection*{Tests destacados}
\begin{itemize}
  \item \textbf{Perft}: validacion de generacion de movimientos en posiciones conocidas
        (startpos, kiwipete, position3) con nodos por profundidad exactos
  \item \textbf{Mate-in-1}: jaque mate por la espalda (Re8\#), jaque mate del erudito (Qxf7\#)
  \item \textbf{Mate-in-2}: patron de Damiano con busqueda a profundidad 6
  \item \textbf{EOF critico}: verifica que el motor no entra en bucle infinito al cerrar stdin
  \item \textbf{Tablero restaurado}: verifica que la busqueda no altera el tablero de entrada
\end{itemize}

% ===========================================================================
\section{Compilacion y ejecucion}
% ===========================================================================

\subsection*{Con Makefile}
\begin{verbatim}
make all          # Compila motor + herramientas neuronales
make engine       # Solo el motor
make neural       # Solo herramientas neuronales
make test         # Compila y ejecuta los 52 tests
make clean        # Limpia binarios
\end{verbatim}

\subsection*{Con CMake (recomendado)}
\begin{verbatim}
cmake -S . -B build -DCMAKE_BUILD_TYPE=Release
cmake --build build -j$(nproc)
cd build && ctest
\end{verbatim}

\subsection*{Jugar}
\begin{verbatim}
./play.sh                          # Profundidad 3
./play.sh --depth 5                # Profundidad 5
./play.sh --fen "FEN"              # Desde una posicion concreta
./play.sh --build                  # Compilar primero
./bin/alphabeta_optimized          # Directamente
\end{verbatim}

\subsection*{Comandos durante la partida}
\begin{verbatim}
e2e4        Mover pieza (notacion UCI)
e7e8q       Promocion de peon
e8g8        Enroque corto
d           Mostrar tablero
fen         Mostrar FEN actual
undo        Deshacer 2 movimientos (tuyo + del motor)
new         Nueva partida
go depth 5  Motor juega con profundidad 5
quit        Salir
\end{verbatim}

% ===========================================================================
\section{Estado del proyecto}
% ===========================================================================

\subsection*{Completado}
\begin{itemize}
  \item Motor con Alpha-Beta + profundizacion iterativa + tabla de transposicion
  \item Quiescence search + ordenacion de movimientos (TT move, MVV-LVA, killers, history)
  \item Evaluacion heuristica (material, peones, seguridad rey, control)
  \item Deteccion correcta de jaque mate, ahogado y tablas
  \item Interfaz con comandos (quit, undo, new, fen, d, go depth)
  \item 52 tests automatizados (unitarios, perft, mate-in-N, end-to-end)
  \item Script \texttt{play.sh} para jugar facilmente
  \item Redes neuronales C++ (junior y grande) con inicializacion He
  \item Entrenamiento PyTorch con mini-batches y exportacion binaria
  \item Perspectiva del lado al mueve en encoding y evaluacion
\end{itemize}

\subsection*{Pendiente}
\begin{itemize}
  \item Podas avanzadas (null-move pruning, LMR, futility)
  \item NeuralEvaluator integrado como \texttt{Evaluator}
  \item Protocolo UCI completo
  \item Multi-hilo (Lazy SMP)
  \item Libro de aperturas
  \item Medicion de fuerza con cutechess-cli
  \item Arquitectura NNUE (acumulador incremental)
\end{itemize}

% ===========================================================================
\section{Licencia}
% ===========================================================================

Este proyecto esta licenciado bajo la
\href{https://creativecommons.org/licenses/by-nc/4.0/}{Creative Commons
Attribution-NonCommercial 4.0 International License} (CC BY-NC 4.0).

\begin{itemize}
  \item \textbf{Permitido}: compartir, adaptar, uso educacional/personal/investigacion
  \item \textbf{No permitido}: uso comercial sin permiso explicito del autor
\end{itemize}

% ===========================================================================
\section{Contacto y agradecimientos}
% ===========================================================================

\textbf{Autor}: Llorenc Perez --- \texttt{perezllorenc@gmail.com}

\textbf{Repositorio}: \url{https://github.com/LLPB-pixel/ajedrezUPV}

Agradecimientos: UPV, Disservin (chess-library), nlohmann (json), Lichess
(base de datos de evaluaciones), PyTorch team.

\vfill
\begin{center}
\textit{Generado: \today}
\end{center}

\end{document}
